\section{Transformer tạo biểu diễn theo ngữ cảnh}
\label{sec:ch02-transformer}

\index{transformer|(}%
\index{cơ chế chú ý}%
Từ sự phân biệt giữa đầu vào, trạng thái và đầu ra ở mục~\ref{sec:ch02-doi-tuong-can-giai-thich},
ta có thể nhìn transformer như cơ
chế biến một dãy token thành các biểu diễn phụ thuộc ngữ cảnh. Bài báo gốc
giới thiệu transformer như một kiến trúc dựa hoàn toàn vào
\tn{cơ chế chú ý}{attention mechanism}, không dùng đệ quy hay tích chập; cấu
trúc encoder-decoder của nó ghép \tn{tự chú ý}{self-attention} với các tầng
truyền thẳng theo từng vị trí~\autocite{p01attention}.

Trong tự chú ý, biểu diễn tại một vị trí tạo truy vấn, khóa và giá trị.
Truy vấn so với các khóa để lấy trọng số, rồi các giá trị được trộn theo những
trọng số ấy. Nhờ vậy, một token có thể lấy thông tin từ các token khác trong
ngữ cảnh thay vì chỉ nhận trạng thái từ vị trí liền trước. Cơ chế chú ý nhiều đầu
thực hiện nhiều phép chiếu song song; mỗi đầu chú ý có thể học một kiểu quan
hệ, dù việc quan sát một đầu có trọng số lớn chưa chứng minh đầu ấy là nguyên
nhân của một token đầu ra.

Điểm cuối cùng này nối kiến trúc với vấn đề độ trung thực. Ma trận attention
là một đại lượng có thể vẽ và so sánh, nên rất dễ bị đọc như một bản đồ của
lý do mô hình trả lời. Nhưng attention mô tả một phép trộn trong một lần tính
toán cụ thể. Muốn kết luận về sự phụ thuộc nhân quả của đầu ra, ta vẫn phải
nêu can thiệp nào được thực hiện và kiểm tra đầu ra đổi ra sao. Chương này vì
thế dùng attention để gọi tên một vị trí trong cơ chế, không dùng nó thay cho
một cách đánh giá.

Các biểu diễn cũng thay đổi qua các tầng. Một token ở tầng đầu có thể giữ
nhiều thông tin bề mặt, trong khi biểu diễn ở tầng sau chịu ảnh hưởng của toàn
bộ ngữ cảnh mà kiến trúc cho phép nó truy cập. Khi một phương pháp hậu nghiệm ở
chương~\ref{ch:lime} hoặc chương~\ref{ch:shap} chỉ nhìn đầu vào và đầu ra,
nó không cần giả định rằng đã biết biểu diễn nào giữ một khái niệm. Đổi lại,
nó phải chịu trách nhiệm về chính can thiệp trên đầu vào mà nó chọn.
\index{transformer|)}%
